\section*{Chapitre \ref{ch:b2:limb-organogenesis} --- L'organogenèse~: la construction du membre des tétrapodes}

\begin{solution}{exo:b2:limb-organogenesis:1}
Proximo-distal~: la \omterm{def:b2:limb-organogenesis:bud}{crête apicale ectodermique} le long de l'extrémité,
les FGF. Antéro-postérieur~: la \omterm{def:b2:limb-organogenesis:bud}{zone d'activité polarisante} au bord
postérieur, \omterm{prop:b2:limb-organogenesis:zpa}{Sonic hedgehog}. Dorso-ventral~: l'ectoderme dorsal, Wnt7a.
\end{solution}

\begin{solution}{exo:b2:limb-organogenesis:2}
Stylopode~: humérus, fémur (Hoxa/d 9--10). Zeugopode~: radius et
cubitus, tibia et péroné (Hox 11). Autopode~: poignet et main, cheville
et pied (Hox 13).
\end{solution}

\begin{solution}{exo:b2:limb-organogenesis:3}
Spécification du champ de membre (Tbx5)~; croissance du bourgeon sous
l'effet du FGF de la crête~; organisation par la ZPA et l'ectoderme
dorsal~; condensation en rayons des cellules qui quittent la \omterm{prop:b2:limb-organogenesis:aer}{zone de
progression}~; différenciation du cartilage~; formation des articulations
là où le cartilage est inhibé~; invasion par les vaisseaux sanguins et
remplacement du cartilage par l'os~; mort des cellules interdigitales~;
migration des précurseurs musculaires issus des somites, qui fusionnent
et s'attachent par des tendons.
\end{solution}

\begin{solution}{exo:b2:limb-organogenesis:4}
(a) Un moignon avec peu ou pas d'humérus~; (b) humérus, radius et cubitus,
mais pas de main~; (c) un membre complet --- le FGF remplace la crête.
\end{solution}

\begin{solution}{exo:b2:limb-organogenesis:5}
$x = 40\ln 2 = \qty{28}{\micro m}$, $40\ln 5 = \qty{64}{\micro m}$,
$40\ln 20 = \qty{120}{\micro m}$~: territoires de 28, 36 et
\qty{56}{\micro m}.
\end{solution}

\begin{solution}{exo:b2:limb-organogenesis:6}
(a) 4-3-2-2-3-4, une série complète en miroir. (b) Une source plus
faible donne un pic plus bas~: seul le doigt de seuil bas est ajouté,
2-2-3-4. (c) Une exposition courte ne spécifie de même que le doigt
surnuméraire le plus antérieur~: 2-2-3-4 --- les cellules intègrent la
concentration au cours du temps.
\end{solution}

\begin{solution}{exo:b2:limb-organogenesis:7}
$5000\times 2^{8} = 1.3\times 10^{6}$. Atteindre $4\times 10^{6}$ exige
$\log_2 800 = 9.6$ doublements~: le cycle doit donc durer en moyenne
environ \qty{10}{h}, ou le décompte doit inclure les précurseurs
musculaires qui migrent depuis les somites~; la mort cellulaire
aggraverait l'écart au lieu de le réduire.
\end{solution}

\begin{solution}{exo:b2:limb-organogenesis:8}
Les cellules interdigitales du poulet reçoivent un signal BMP et
meurent~; celles du canard expriment un inhibiteur de BMP et survivent,
si bien que la palmure persiste. L'inhibiteur apporté sur une bille
préserve la palmure du poulet~: une patte palmée. Une mort à l'heure
dans une greffe signifie que les cellules avaient déjà reçu leurs
instructions avant le prélèvement --- la décision est prise tôt et
exécutée plus tard, de manière autonome.
\end{solution}

\begin{solution}{exo:b2:limb-organogenesis:9}
Sans Hoxa13 ni Hoxd13~: pas d'autopode --- le membre antérieur s'arrête
au poignet. Sans Hoxa11 ni Hoxd11~: un radius et un cubitus très
raccourcis, avec une main fixée presque directement à l'humérus. Chaque
segment exige sa propre paire de gènes Hox~; les domaines ne sont pas de
simples marqueurs, mais des instructions.
\end{solution}

\begin{solution}{exo:b2:limb-organogenesis:10}
Le FGF de la crête maintient l'expression de Shh dans la ZPA~; Shh, par
un relais mésodermique, maintient l'expression de FGF dans la crête. Le
signal de chaque centre certifie donc que l'autre fonctionne~: un membre
ne croît jamais sans être organisé ni n'est organisé sans croître, et
une perte transitoire de l'un est réparée par l'autre. Si Shh ne pouvait
pas maintenir le FGF, la crête régresserait en une journée et la
croissance s'arrêterait au segment atteint~: un membre tronqué. La
boucle du travail de l'accouchement prend fin avec l'événement qu'elle
provoque --- l'expulsion supprime le stimulus~; celle du membre prend
fin avec la différenciation --- le mésoderme cesse de relayer lorsque le
membre est achevé.
\end{solution}

\begin{solution}{exo:b2:limb-organogenesis:11}
$e^{4r} = 8$~: $r = \ln 8/4 = \qty{0.52}{d^{-1}}$. Croissance~: une
prolifération plus longue et plus rapide du cartilage des doigts (un
amplificateur du membre qui augmente l'expression d'un gène de
croissance dans les doigts). Mort~: la survie du tissu interdigital
grâce à l'expression d'un inhibiteur de BMP, qui conserve la membrane.
\end{solution}

\begin{solution}{exo:b2:limb-organogenesis:12}
Les signaux sont réutilisés parce qu'un module récepteur--transduction
fonctionnel est peu coûteux à redéployer~: ce qui change d'un tissu à
l'autre, ce n'est pas le signal mais l'ensemble des gènes que les
cellules réceptrices sont compétentes pour activer, si bien que le même
profil de concentration signifie motoneurone à un endroit et doigt 4 à
un autre. Une \omterm{def:b2:genome-diversification:mutation}{mutation} d'un tel gène touche donc plusieurs organes à la
fois --- les \omterm{def:b2:genome-diversification:mutation}{mutations} de Shh associent holoprosencéphalie et anomalies
des membres --- à moins qu'elle ne porte sur un amplificateur propre à un
tissu, auquel cas elle ne modifie qu'un seul organe, et c'est ainsi que
se produit l'essentiel du changement évolutif de ces gènes.
\end{solution}

\begin{solution}{pb:b2:limb-organogenesis:1}
\textbf{1.} $96/12 = 8$ doublements~: $2000\times 256 = 5.1\times 10^{5}$
cellules.
\textbf{2.} $1000/256 \approx 4$~: agrandissement des cellules et
matrice cartilagineuse.
\textbf{3.} Longueur $\times 10$~; un cylindre de rayon constant
donnerait $\times 10$ en volume~; le facteur cent supplémentaire signifie
que le rayon a lui aussi été multiplié par dix environ --- le bourgeon
s'est élargi en palette en s'allongeant.
\textbf{4.} Au 3\textsuperscript{e} jour, $300/400 = \qty{75}{\%}$~; au
7\textsuperscript{e} jour, $300/4000 =
\qty{7.5}{\%}$.
\textbf{5.} Seule l'extrémité est sous l'effet du signal~; à mesure que
l'extrémité avance, les cellules laissées en arrière quittent la zone
dans l'ordre où elles ont été produites, si bien que les structures
proximales sont spécifiées en premier et les distales en dernier.
\textbf{6.} Le zeugopode (radius et cubitus), spécifié entre 3.5 et
4 jours.
\textbf{7.} Stylopode à 3.5 jours, zeugopode à 4.0 jours, autopode à
4.5 jours (extrémités des doigts plus tard).
\textbf{8.} Une demi-journée~: un cycle cellulaire.
\textbf{9.} Chaque cycle passé dans la \omterm{prop:b2:limb-organogenesis:aer}{zone de progression} fait avancer
le code Hox (9--10, puis 11, puis 13)~; les cellules qui partent après
un, deux ou trois cycles portent le code du stylopode, du zeugopode ou de
l'autopode.
\textbf{10.} Une aile complète~: le FGF est ce que fournissait la crête.
\textbf{11.} Une seconde excroissance sur la face dorsale~: des
structures distales dupliquées, un membre surnuméraire.
\textbf{12.} Le mésoderme cesse de relayer Shh vers la crête,
l'expression de Shh s'arrête, la crête régresse, et les cellules se
différencient au lieu de se diviser~: la boucle qui entretenait la
croissance est rompue.
\textbf{13.} $k = \ln 2/1740 = \qty{4.0e-4}{s^{-1}}$~; $\lambda =
\sqrt{10^{-12}/4\times 10^{-4}} = \qty{50}{\micro m}$.
\textbf{14.} $x = 50\ln(1/0.6) = 26$, $50\ln 4 = 69$, $50\ln 12.5 =
\qty{126}{\micro m}$~: territoires de 26, 43 et \qty{57}{\micro m}.
\textbf{15.} $600 - 126 = \qty{474}{\micro m}$.
\textbf{16.} Chaque frontière s'éloigne de $50\ln 2 = \qty{35}{\micro
m}$~: 61, 104, 161~; les territoires des doigts 3 et 2 gardent leur
largeur, celui du doigt 4 passe de 26 à \qty{61}{\micro m}, et la zone
organisée s'étend de 126 à \qty{161}{\micro m}, assez pour un doigt 2
antérieur surnuméraire.
\textbf{17.} Chaque série exige \qty{126}{\micro m}~: au moins
\qty{252}{\micro m}~; le bourgeon de \qty{600}{\micro m} a la place
nécessaire, avec \qty{348}{\micro m} de tissu sans doigt au milieu.
\textbf{18.} À \qty{200}{\micro m}, les deux gradients se chevauchent et
s'additionnent~: les cellules du milieu reçoivent assez de signal pour
former un doigt 3 et les territoires du doigt 2 disparaissent --- 4-3-4
ou 4-3-3-4, moins de doigts qu'une série complète en miroir.
\textbf{19.} Un quart des cellules donne $c_0/4$~: aucune cellule
n'atteint le seuil du doigt 4 ou du doigt 3 ($0.25 c_0$ au bord
lui-même), et le seuil du doigt 2 est franchi jusqu'à
$50\ln(0.25/0.08) =
\qty{57}{\micro m}$~: un seul doigt 2 surnuméraire.
\textbf{20.} $3\times 300\times 300\times 200 = 5.4\times 10^{7}$
\unit{\micro m^3}~: $5.4\times 10^{4}$ cellules~; sur \qty{86400}{s},
$0.6$ cellule par seconde.
\textbf{21.} Bille de BMP dans la palmure du canard~: la palmure meurt
et les doigts se séparent~; inhibiteur dans la palmure du poulet~: la
palmure survit, d'où une patte palmée.
\textbf{22.} $e^{0.5\times 4} = e^{2} = 7.4$. La seconde modification est
la survie de la membrane interdigitale grâce à un inhibiteur de BMP.
\textbf{23.} La phase tardive d'expression de Hox 13 qui spécifie un
autopode, et la persistance de la crête~: le repli du poisson forme à la
place des rayons de nageoire. La nageoire de \emph{Tiktaalik} contient
des os semblables à ceux d'un poignet --- une nageoire qui commence à
devenir un membre.
\textbf{24.} $5.4\times 10^{4}/1000 = 54$~: $\log_2 54 = 5.8$ doublements,
environ \qty{69}{h}, trois jours pour fabriquer ce qu'une journée
élimine.
\textbf{25.} Huit doublements~; frontières des doigts à 26, 69 et
\qty{126}{\micro m}~; une duplication en miroir exige au moins
\qty{252}{\micro m}.
\end{solution}
